\chapter{て形② — Forma て (II): Permisos y Peticiones}
\unidadheader{14}{て形②}{Forma て (II)}{Pedir favores, dar permisos y expresar prohibiciones}

\begin{tcolorbox}[vocabbox, title={📌 Vocabulario Nuevo}]
\begin{tabularx}{\linewidth}{llX}
\toprule
\textbf{Japonés} & \textbf{Español} & \textbf{Tipo} \\
\midrule
\vocabitem{使う}{つかう}{usar}{v. gr.1}
\vocabitem{置く}{おく}{poner / colocar}{v. gr.1}
\vocabitem{捨てる}{すてる}{tirar / desechar}{v. gr.2}
\vocabitem{手伝う}{てつだう}{ayudar}{v. gr.1}
\vocabitem{急ぐ}{いそぐ}{apurarse}{v. gr.1}
\vocabitem{休む}{やすむ}{descansar / faltar}{v. gr.1}
\vocabitem{続ける}{つづける}{continuar}{v. gr.2}
\vocabitem{始める}{はじめる}{comenzar (algo)}{v. gr.2}
\vocabitem{気をつける}{きをつける}{tener cuidado}{v. gr.2}
\vocabitem{片付ける}{かたづける}{ordenar / limpiar}{v. gr.2}
\bottomrule
\end{tabularx}
\end{tcolorbox}

\subsection*{🗣️ Frases Clave}

\frase{ここに\ruby{座}{すわ}ってもいいですか？}{¿Puedo sentarme aquí?}
\frase{どうぞ。}{Adelante. / Por favor.}
\frase{ここでたばこを\ruby{吸}{す}ってはいけません。}
  {Aquí no está permitido fumar.}
\frase{\ruby{窓}{まど}を\ruby{開}{あ}けてください。}{Abra la ventana, por favor.}

\subsection*{⚙️ Gramática}

\patron{Petición amable}
  {V-て + ください}
  {\ej{\ruby{名前}{なまえ}を\ruby{書}{か}いてください。}
    {Por favor, escriba su nombre.}
  \ej{もう\ruby{少}{すこ}しゆっくり\ruby{話}{はな}してください。}
    {Por favor, hable un poco más despacio.}}

\patron{Pedir permiso}
  {V-て + もいいですか？}
  {\ej{\ruby{写真}{しゃしん}を\ruby{撮}{と}ってもいいですか？}
    {¿Puedo tomar una foto?}
  \ej{トイレを\ruby{使}{つか}ってもいいですか？}
    {¿Puedo usar el baño?}}

\patron{Prohibición}
  {V-て + はいけません / はだめです}
  {\ej{ここで\ruby{走}{はし}ってはいけません。}{No está permitido correr aquí.}
  \ej{ゴミを\ruby{捨}{す}ててはだめです。}{No debes tirar basura.}}

\begin{tcolorbox}[ejerciciobox, title={📝 Ejercicios Rápidos}]
\begin{enumerate}
  \item Pide a alguien que espere (待つ) en japonés.
  \item Pregunta si puedes usar un diccionario (辞書).
  \item Traduce la prohibición: No está permitido entrar aquí.
\end{enumerate}
\vspace{0.4em}
\textbf{Respuestas:}
1) \ruby{待}{ま}ってください。\quad
2) \ruby{辞書}{じしょ}を\ruby{使}{つか}ってもいいですか？\quad
3) ここに\ruby{入}{はい}ってはいけません。
\end{tcolorbox}

\begin{tcolorbox}[kanjibox, title={🀄 Kanjis de la Unidad}]
\begin{tabularx}{\linewidth}{cllXX}
\toprule
\textbf{Kanji} & \textbf{On} & \textbf{Kun} & \textbf{Significado} & \textbf{Vocab} \\
\midrule
\kanjirow{使}{し}{つか}{usar}{\ruby{使}{つか}う / \ruby{使用}{しよう}}
\kanjirow{置}{ち}{お}{poner / dejar}{\ruby{置}{お}く / \ruby{位置}{いち}}
\kanjirow{急}{きゅう}{いそ}{urgente / apurarse}{\ruby{急}{いそ}ぐ / \ruby{急行}{きゅうこう}}
\kanjirow{休}{きゅう}{やす}{descansar}{\ruby{休}{やす}む / \ruby{休日}{きゅうじつ}}
\kanjirow{始}{し}{はじ}{comenzar}{\ruby{始}{はじ}める / \ruby{開始}{かいし}}
\bottomrule
\end{tabularx}
\end{tcolorbox}
